\documentclass[11pt]{article}

% Encoding
\usepackage[utf8]{inputenc}
\usepackage[T1]{fontenc}

% Professional preprint style: fonts, geometry, headings, title, header.
\usepackage{arxiv}

% Content packages
\usepackage{amsmath}
\usepackage{amssymb}
\usepackage{graphicx}
\usepackage{booktabs}
\usepackage{listings}
\usepackage{subcaption}
\usepackage{enumitem}
\usepackage{natbib}

% hyperref loaded late, then coloured to match the palette.
\usepackage{hyperref}
\usepackage{url}
\hypersetup{
  colorlinks=true,
  linkcolor=arxivaccent,
  citecolor=arxivaccent,
  urlcolor=arxivaccent,
  filecolor=arxivaccent,
  breaklinks=true,
}

% Python code listing style
\definecolor{codebg}{rgb}{0.95,0.95,0.97}
\definecolor{codegreen}{rgb}{0.0,0.5,0.0}
\definecolor{codegray}{rgb}{0.5,0.5,0.5}
\definecolor{codepurple}{rgb}{0.58,0,0.82}
\definecolor{codeblue}{rgb}{0.0,0.0,0.7}

\lstdefinestyle{pythonstyle}{
    backgroundcolor=\color{codebg},
    basicstyle=\ttfamily\small,
    breakatwhitespace=false,
    breaklines=true,
    captionpos=b,
    commentstyle=\color{codegreen},
    keywordstyle=\color{codeblue}\bfseries,
    numberstyle=\tiny\color{codegray},
    stringstyle=\color{codepurple},
    language=Python,
    showstringspaces=false,
    numbers=left,
    numbersep=5pt,
    frame=single,
    rulecolor=\color{codegray},
    tabsize=4,
    morekeywords={dataclass, List, Optional, date, assert, Dict, field}
}

\lstdefinestyle{treestyle}{
    basicstyle=\ttfamily\small,
    frame=none,
    numbers=none,
    breaklines=false,
    backgroundcolor=\color{codebg},
}

% Style for conversation transcript excerpts (italic dialogue, framed).
\definecolor{convbg}{rgb}{0.98,0.97,0.92}
\definecolor{speakerblue}{rgb}{0.10,0.30,0.60}
\lstdefinestyle{conversationstyle}{
    backgroundcolor=\color{convbg},
    basicstyle=\ttfamily\footnotesize,
    breakatwhitespace=false,
    breaklines=true,
    captionpos=b,
    language={},          % plain transcript: no Python string/keyword colouring
    keywordstyle={},
    stringstyle={},
    commentstyle={},
    showstringspaces=false,
    numbers=none,
    frame=single,
    rulecolor=\color{codegray},
    tabsize=2,
}

% Style for runtime output / alert blocks (red border, fixed-width).
\definecolor{alertbg}{rgb}{1.00,0.96,0.94}
\definecolor{alertred}{rgb}{0.70,0.10,0.10}
\lstdefinestyle{alertstyle}{
    backgroundcolor=\color{alertbg},
    basicstyle=\ttfamily\footnotesize,
    breakatwhitespace=false,
    breaklines=true,
    captionpos=b,
    commentstyle=\color{codegreen},
    keywordstyle=\color{alertred}\bfseries,
    showstringspaces=false,
    numbers=none,
    frame=single,
    rulecolor=\color{alertred},
    tabsize=2,
    morekeywords={CRITICAL,WARNING,INFO,ALERT,DRUG-ALLERGY,CONFLICT},
}

\lstset{style=pythonstyle}


\title{User as Code: Executable Memory for Personalized Agents}

\author{
  Bojie Li \\
  Pine AI
}

\date{}
\runningtitle{User as Code: Executable Memory for Personalized Agents}

\begin{document}
\maketitle

\begin{abstract}
A personalized AI agent needs a \emph{user memory}: a persistent model of who the user is, accumulated across many conversations and consulted on each new one. Today this memory is almost always stored as unstructured text, a knowledge graph, or a flat store of extracted facts, and consulted by \emph{retrieval}---fetching the stored entries most similar to the current request. Such ``bag-of-facts'' memory recalls individual facts well, but because storing a fact and acting on it are separate steps, it struggles to resolve contradictions, aggregate over many records, or enforce logical rules. We argue that user memory should instead be \emph{executable}. We introduce \textbf{User as Code} (UaC), a paradigm in which an agent's model of a user is a living software project: typed Python objects hold the user's state, and ordinary Python functions encode the rules that govern it, so that representing the user and reasoning about the user happen in one medium an interpreter can run. The enabling mechanism is a \emph{two-phase pipeline} that turns raw conversation into this code---an append-only log that never discards a fact, periodically checkpointed into structured, typed code---a design long used in database systems and, to our knowledge, applied here for the first time to LLM memory.

This shift changes what memory can do. On standard long-term conversation benchmarks, UaC is competitive with both a full-context upper bound and the strongest prior memory systems on ordinary factual recall (78.8\% on LOCOMO). Its advantage emerges where representation matters most. On aggregate questions over a user's history---\emph{``how many international trips did I take last year?''}---retrieval-based memory collapses (6--43\%) while UaC stays near-perfect (99\%), because the answer is a one-line computation over typed state rather than a search over text. And because its rules, once written, execute deterministically whenever the state changes, UaC can surface unsolicited, safety-critical alerts---such as a newly prescribed drug that conflicts with an allergy recorded months earlier---a proactive capability that query-driven memory cannot provide at all. The one-time cost of structuring the code is repaid after only a few queries against the same user.
\end{abstract}

\begin{center}
\small
Code: \url{https://github.com/19PINE-AI/user-as-code} \\[2pt]
Website: \url{https://01.me/research/user-as-code}
\end{center}
\vspace{-0.6em}

\section{Introduction}

\subsection{The Memory Challenge for Personalized Agents}

To build truly personalized, continuous-service AI agents, a \emph{User Memory} system is indispensable. Memory is not simply recording every word a user says---it is gradually forming a model of the user through sustained interaction: their personality, habits, values, and circumstances. Agent memory spans multiple layers: working memory (single-session trajectory), long-term memory (persistent cross-session knowledge), and business state (task-phase abstractions). This paper focuses on long-term user memory.

Memory consists of two fundamental components---\textbf{memorizing} (extracting structured knowledge from raw interactions and persisting it) and \textbf{retrieval} (recalling the right knowledge at the right time). All memory systems must address both halves. User as Code contributes to both: on the memorizing side, knowledge is extracted into structured, typed, version-controlled code; on the retrieval side, executing code against the state is itself a form of retrieval---the generate-verify-review loop \emph{computes} answers that no similarity search could produce.

\subsection{Three Categories of Memory Tasks}

User memory capabilities can be understood along a spectrum of increasing difficulty:

\begin{itemize}[leftmargin=*]
\item \textbf{Basic Recall:} Store and retrieve directly provided, unambiguous information (e.g., ``My membership number is 12345''). All existing approaches handle this adequately.

\item \textbf{Multi-session Retrieval:} Retrieve and reason over information scattered across sessions, entities, and time---requiring disambiguation (``which car?''), compound event resolution (``cancel the LA trip'' = flights + hotels + car rentals), and conflict resolution (multi-party conflicting instructions).

\item \textbf{Active Service:} Proactively synthesize information across many sessions to provide anticipatory, unsolicited help---e.g., cross-referencing a new flight with passport info stored months ago to alert that the passport expires too soon. No retrieval sophistication can reliably compute that a passport expiring Feb~18 leaves only 34 days of validity for a Jan~15 departure---far short of the 180-day requirement. This is \emph{arithmetic}, not a similarity match.
\end{itemize}

Existing benchmarks (LOCOMO~\cite{maharana2024locomo}, LongMemEval~\cite{wu2025longmemeval}, MemoryAgentBench~\cite{hu2025memoryagentbench}) primarily test the first two categories. LoCoMo-Plus~\cite{li2026locomoplus} tests implicit constraints and MemoryArena~\cite{he2026memoryarena} shows that agents saturated on LOCOMO still fail in agentic settings---but no benchmark tests proactive memory-triggered alerting. User as Code is specifically designed to enable Active Service, while remaining competitive on the first two categories.

\subsection{The ``Bag of Facts'' Problem}

Current systems---vector databases, JSON stores, knowledge graphs, flat fact extractors---treat memory as a loose collection of facts. This creates two fundamental problems:

\textbf{Conflict:} ``I love cilantro'' vs.\ later ``I actually hate cilantro''---vector retrieval might return both or the wrong one. Without a single source of truth with version history, conflicting facts coexist unresolved.

\textbf{Logic Gap:} Even the most advanced structured approaches---from atomic fact notes through JSON cards to knowledge graphs---cannot express rules like ``if passport expires within 180 days of travel date, raise an alert.'' This requires \emph{executable} logic that the system can \textbf{generate on the fly} as novel situations arise.

\textbf{Thesis:} The fundamental limitation of all existing formats---natural language, JSON, knowledge graphs---is that they separate \emph{representation} from \emph{verification}. Code is the only format where the data and the checks over that data coexist in the same medium. User as Code exploits this unification.

\subsection{The User as Code Proposal}

We model a user's entire memory as a \textbf{self-evolving software project}---with domain packages, schema definitions, state data, cross-domain constraints, and tests. The core insight is that the same Python that stores \texttt{passport\_expiry = date(2025, 2, 18)} can also express---and \emph{execute}---\texttt{assert (passport\_expiry - flight\_date).days >= 180}. No other format unifies storage and verification this way.

Modern AI agents (Manus, Cursor, Claude Code, Windsurf) are fundamentally coding agents built on virtual file systems. User as Code leverages this: the user project is a \textbf{directory in the agent's own workspace}, manipulated with the same file and interpreter primitives the agent already uses. No custom memory API is needed.

The \textbf{generate-verify-review loop} is a continuous cycle: the coding agent generates Python checks, the interpreter verifies them deterministically, and the agent reviews results to refine or persist. This directly \textbf{reduces hallucination} by delegating computation to deterministic execution rather than unreliable natural-language arithmetic.

% ===================================================================
\section{Related Work}
% ===================================================================

The field of LLM agent memory has experienced explosive growth in 2025--2026, with 20+ memory systems, 15+ benchmarks, multiple comprehensive surveys~\cite{hu2025memoryage, du2026memorysurvey, yang2026graphmemory, wu2025humantoai, zhang2024memorysurvey}, and a dedicated ICLR 2026 workshop (MemAgents). We organize related work around three gaps that motivate User as Code.

\subsection{Gap 1: Representation Separates Storage from Verification}

All existing memory systems store user state in formats that cannot be directly executed.

\textbf{Flat fact and note-based systems} extract atomic facts from conversations. Mem0~\cite{chhikara2025mem0} uses CRUD lifecycle with LLM-based merge/update, achieving 26\% improvement over OpenAI's memory on LOCOMO. A-MEM~\cite{xu2025amem} applies the Zettelkasten method---structured notes with dynamic linking. ENGRAM~\cite{patel2025engram} shows that careful memory typing with dense retrieval matches complex architectures at $\sim$1\% token cost. These systems produce readable, retrievable memory---but facts remain inert.

\textbf{Knowledge graph approaches} add relational structure. Zep/Graphiti~\cite{rasmussen2025zep} builds temporal KGs with bitemporal modeling. MAGMA~\cite{jiang2026magma} uses multi-graph traversal across semantic, temporal, causal, and entity dimensions. Memoria~\cite{sarin2025memoria} combines session summarization with weighted KGs. These excel at entity relations but conditional logic degrades when forced into triples.

\textbf{Hybrid architectures} combine multiple strategies. Hindsight~\cite{latimer2025hindsight} maintains four logical networks achieving 91.4\% on LongMemEval. MemMachine~\cite{wang2026memmachine} stores entire episodes, achieving 0.9169 on LOCOMO (current SOTA). D-Mem~\cite{you2026dmem} applies dual-process theory. These push recall high but still represent state as text.

\textbf{OS-inspired systems} use programmatic infrastructure without programmatic representation. MemGPT/Letta~\cite{packer2023memgpt} treats the LLM as an OS managing tiered memory. MemOS~\cite{li2025memos} introduces MemCube with OS-style scheduling. LangMem~\cite{langmem2025} allows agents to update their own system instructions---the closest approach to ``executable memory.'' But in all cases, user state remains text. \textbf{User as Code closes this gap: code is not just the infrastructure for managing memory, but the representation of memory itself.}

\subsection{Gap 2: No Benchmark Tests Proactive Memory-Triggered Alerting}

\textbf{Standard benchmarks} test reactive recall. LOCOMO~\cite{maharana2024locomo} provides 1,986 QA pairs; LongMemEval~\cite{wu2025longmemeval} tests 500 questions across five abilities. \textbf{Cognitive benchmarks} push further: LoCoMo-Plus~\cite{li2026locomoplus} tests implicit constraints under cue-trigger semantic disconnect; MemoryArena~\cite{he2026memoryarena} shows LOCOMO-saturated agents fail in agentic settings; MemoryAgentBench~\cite{hu2025memoryagentbench} tests selective forgetting. ProAgentBench~\cite{tang2026proagentbench} tests proactive task assistance but not memory-triggered alerting.

PersistBench~\cite{pulipaka2026persistbench} identifies memory safety risks: cross-domain leakage (53\% failure) and memory-induced sycophancy (97\% failure). HaluMem~\cite{chen2025halumem} benchmarks hallucination in memory operations. MemoryBench~\cite{ai2025memorybench} finds no advanced system consistently outperforms simple RAG.

\textbf{No existing benchmark tests the initiation asymmetry} central to Active Service: the system must generate alerts \emph{without being asked}, triggered by state changes rather than queries.

\subsection{Gap 3: Memory Operations Are Not Yet Optimized End-to-End}

Recent work demonstrates RL-trained memory management. Memory-R1~\cite{yan2025memoryr1} trains ADD/UPDATE/DELETE policies with GRPO, improving F1 by 48\%. AgeMem~\cite{yu2026agemem} learns proactive summarization through progressive RL. Mem-alpha~\cite{wang2025memalpha} shows generalization to 400K+ tokens despite training on 30K. MemFactory~\cite{guo2026memfactory} provides the first unified training+inference framework.

Cognitively-inspired systems ground memory in neuroscience: FadeMem~\cite{wei2026fademem} implements biologically-inspired forgetting; LightMem~\cite{fang2025lightmem} follows the Atkinson-Shiffrin model; EverMemOS~\cite{hu2026evermemos} uses engram-inspired lifecycles (93.05\% on LOCOMO).

All these train over custom memory APIs. \textbf{User as Code's file-system-based architecture makes memory operations native file actions}---directly compatible with RL training without custom toolset wrappers.

\subsection{Commercial Systems}

Claude Memory (Anthropic, 2025) uses markdown files in a hierarchical structure---the closest commercial approach to our file-system philosophy, though not executable. ChatGPT Memory (OpenAI, 2025) maintains a four-component architecture. ChatGPT Pulse is the only deployed proactive system, synthesizing memory overnight for morning briefings.

% ===================================================================
\section{Methodology: The User as Code Architecture}
% ===================================================================

A user is a \textbf{self-evolving software project} that lives in the agent's own workspace. The agent reads, writes, and executes this project using the same file system and interpreter tools it already has.

\subsection{Design Principles}

\begin{enumerate}[leftmargin=*]
\item \textbf{Separation of Structure and Data.} Schema (class definitions) is separated from state (instance data) and archive (historical records).
\item \textbf{Modularity by Life Domain.} Memory is partitioned into independent domain packages (\texttt{travel/}, \texttt{finance/}, \texttt{health/}), each loadable and testable independently.
\item \textbf{Progressive Disclosure.} The agent navigates via a compact manifest (always loaded, $\sim$200--300 tokens) and retrieves domain modules on demand.
\item \textbf{Unified Representation and Verification.} Code serves as both the storage format the LLM reads and the verification medium the interpreter executes.
\item \textbf{Agent-Native File System.} The user project is a directory in the agent's workspace---no custom memory API needed.
\end{enumerate}

\subsection{The User Project Structure}

\begin{lstlisting}[style=treestyle,caption={User project directory structure. One user = one project.},label=lst:structure]
jessica_thompson/
|-- manifest.py          # Always in agent context
|-- domains/
|   |-- identity/
|   |   |-- schema.py    # PersonalInfo dataclass
|   |   +-- state.py     # Current identity data
|   |-- travel/
|   |   |-- schema.py    # Trip, PassportInfo dataclasses
|   |   +-- state.py     # Passport, trips, preferences
|   |-- health/
|   |   |-- schema.py    # Allergy, Medication dataclasses
|   |   +-- state.py     # Allergies, current Rx
|   +-- finance/ ...
|-- constraints/          # LLM-generated, persistent
|   |-- travel_readiness.py
|   +-- health_safety.py
+-- tests/                # LLM-generated invariants
    +-- test_travel.py
\end{lstlisting}

Domain modules are bounded ($\sim$10--20); \texttt{state.py} holds only current active state; history lives in the archive. Domains are added, refactored, and retired---not hardcoded. For a new user, the agent bootstraps the project from the first conversation.

\subsection{Memory Organization}

Memory is organized by \emph{type} (how the LLM interacts with it) and by \emph{tier} (access pattern).

\begin{table}[t]
\caption{Five memory types in User as Code.}
\label{tab:memory-types}
\small
\begin{tabular}{@{}llp{3.5cm}@{}}
\toprule
\textbf{Type} & \textbf{Interaction} & \textbf{Example} \\
\midrule
Factual State & Read & \texttt{expiry = date(2025,2,18)} \\
Preferences & Read & \texttt{"Aisle on long flights"} \\
Persistent Constraints & Auto-Execute & \texttt{travel\_readiness.check()} \\
Ad-hoc Verification & Generate+Execute & \texttt{print((expiry-dep).days)} \\
Narrative Context & Search (RAG) & \texttt{archive.search("JAL")} \\
\bottomrule
\end{tabular}
\end{table}

\textbf{Three tiers:} Tier~1 (Schema---class definitions, bounded), Tier~2 (State---current instance data, moderate), Tier~3 (Archive---full history, unbounded, accessed via RAG).

\subsection{Progressive Disclosure via the Manifest}

The manifest is always in the agent's context ($\sim$200--300 tokens):

\begin{lstlisting}[caption={manifest.py --- always loaded.},label=lst:manifest]
"""User: Jessica Thompson | Updated: 2025-02-14"""

DOMAINS = {
  "travel":  "3 trips | passport exp 2025-02-18",
  "health":  "Allergies: peanuts, penicillin",
  "finance": "Pending wire $15k (CONFLICTING)",
}

ACTIVE_ALERTS = [
  "[CRITICAL] Passport only 34 days valid "
  "for Tokyo trip (need 180)",
  "[CRITICAL] Amoxicillin contraindicated "
  "-- penicillin allergy",
]
\end{lstlisting}

Disclosure levels: L0 (manifest, always) $\rightarrow$ L1 (schema, on topic) $\rightarrow$ L2 (state, on need) $\rightarrow$ L3 (archive, via RAG tool).

\subsection{The Generate-Verify-Review Loop}

The core mechanism for reducing LLM hallucination:

\begin{enumerate}[leftmargin=*]
\item \textbf{Generate:} The agent writes Python constraints against the code-represented state.
\item \textbf{Verify:} Execute in a sandbox---results are deterministic.
\item \textbf{Review:} The agent reviews results, decides to refine, persist, or incorporate.
\end{enumerate}

The agent handles \emph{semantic reasoning} (what to check); the interpreter handles \emph{computation} (correctness). When an ad-hoc check proves useful, the agent promotes it to \texttt{constraints/} as a persistent function. Alerts surface in \texttt{ACTIVE\_ALERTS}, visible at the start of every future conversation. Constraints are not pre-authored---they emerge from the LLM's own reasoning.

\subsection{The Write Path}

Memory separates into session-scoped working directory and cross-session persistent storage. The update pipeline: (1)~classify affected domains, (2)~load schema and state, (3)~diff against existing state, (4)~patch the state file, (5)~validate via tests and constraints, (6)~commit or reject, (7)~regenerate manifest. Every patch includes a source-session reference for audit.

\subsection{End-to-End Walkthrough}

\begin{figure*}[t]
\centering
\small
\begin{tabular}{@{}p{0.23\textwidth}p{0.23\textwidth}p{0.23\textwidth}p{0.23\textwidth}@{}}
\toprule
\textbf{Session 1 (Oct 2024)} & \textbf{Session 8 (Dec 2024)} & \textbf{Session 9 (Dec 2024)} & \textbf{Session 10 (Jan 2025)} \\
\midrule
User mentions passport renewal, expiry Feb~18, 2025. &
User books Tokyo trip, departing Jan~15. &
Agent generates ad-hoc constraint; interpreter outputs ``34 days < 180.'' Agent promotes to persistent constraint. &
User starts unrelated conversation. Agent reads manifest, sees \texttt{ACTIVE\_ALERTS}, proactively warns about passport. \\
\midrule
\multicolumn{4}{l}{\textbf{What baselines produce:} Mem0 stores two disconnected facts. Zep links passport--trip in KG but cannot compute ``34 < 180.''} \\
\multicolumn{4}{l}{\textbf{No baseline produces a proactive alert.}} \\
\bottomrule
\end{tabular}
\caption{End-to-end walkthrough: from conversation to proactive alert across four sessions.}
\label{fig:walkthrough}
\end{figure*}


% ===================================================================
\section{Experiments and Evaluation}
% ===================================================================

\subsection{Experimental Setup}

\textbf{Benchmarks.} We evaluate on three established benchmarks and one novel evaluation protocol:
\begin{itemize}[leftmargin=*]
\item \textbf{LOCOMO}~\cite{maharana2024locomo}: 10 conversations, $\sim$300 turns each, 1,986 QA pairs across single-hop, multi-hop, temporal, and open-domain categories.
\item \textbf{LongMemEval}~\cite{wu2025longmemeval}: 500 curated questions testing five memory abilities at 115K-token scale.
\item \textbf{LoCoMo-Plus}~\cite{li2026locomoplus}: Tests cognitive memory---whether systems can act on implicit constraints under cue-trigger semantic disconnect.
\item \textbf{Active Service protocol} (novel): 40 scenarios across 5 constraint categories (travel document validity, drug interactions, financial authorization conflicts, scheduling conflicts, warranty/deadline expirations). Each scenario seeds facts across 2--4 sessions, then measures whether the system generates unsolicited alerts. The key distinction from all existing benchmarks: \emph{who initiates}---there is no query, only a state change.
\end{itemize}

\textbf{Baselines.} Five systems spanning major architectural approaches: Mem0~\cite{chhikara2025mem0} (flat facts), A-MEM~\cite{xu2025amem} (structured notes), Zep/Graphiti~\cite{rasmussen2025zep} (temporal KG), Hindsight~\cite{latimer2025hindsight} (multi-network), ENGRAM~\cite{patel2025engram} (typed retrieval). All configured following published settings under the same foundation model.

\subsection{Standard Recall and Retrieval}

\textbf{Metrics:} F1, BLEU-1, LLM-as-Judge accuracy, token consumption.

\textbf{Hypothesis:} User as Code achieves competitive performance on standard recall tasks. The typed schema and progressive disclosure are at least as effective as vector retrieval for fact lookup, and version-controlled state provides natural conflict resolution.

\subsection{Cognitive Memory and Implicit Constraints}

\textbf{Metrics:} Constraint consistency scores with evidence-grounded LLM judges.

\textbf{Hypothesis:} User as Code outperforms baselines because code forces ontological commitment---entities must be explicitly resolved at write time---and enables executable verification of latent constraints.

\subsection{Active Service}

The key distinction from all existing benchmarks: there is no query. The trigger is a state change, and the system must (1)~notice that facts in different domains interact, (2)~compute a constraint, (3)~decide this warrants alerting, and (4)~surface the alert unsolicited.

We test two sub-conditions: \textbf{Seen Constraints} (pre-computed alerts in manifest) and \textbf{Unseen Constraints} (agent must generate ad-hoc checks). All baselines fundamentally fail because (a)~without a query, no retrieval is triggered, and (b)~even with co-located facts, computing constraints requires arithmetic, not similarity matching.

\textbf{Preliminary results} using Gemini 3 Flash across 40 scenarios (Table~\ref{tab:active-service}) confirm the hypothesis. We test four conditions: (1)~\emph{Code+Alerts}: full User as Code pipeline with pre-computed constraint alerts in the manifest; (2)~\emph{Code-only}: code representation without pre-computed alerts; (3)~\emph{Flat}: Mem0-style fact list; (4)~\emph{None}: no memory. All conditions use the same neutral system prompt (``you are a personal assistant'')---no instruction to check for issues.

\begin{table}[t]
\caption{Active Service evaluation: proactive alert detection rate across 40 scenarios. Code+Alerts = full User as Code pipeline.}
\label{tab:active-service}
\small
\begin{tabular}{@{}lcccc@{}}
\toprule
\textbf{Category} & \textbf{Code+Alerts} & \textbf{Code} & \textbf{Flat} & \textbf{None} \\
\midrule
Travel documents & \textbf{8/8} & 6/8 & 5/8 & 0/8 \\
Health/drug safety & 7/8 & 7/8 & 7/8 & 0/8 \\
Financial conflicts & 6/8 & 6/8 & 6/8 & 0/8 \\
Scheduling & 7/8 & 6/8 & 7/8 & 0/8 \\
Warranty/deadline & \textbf{7/8} & 4/8 & 4/8 & 0/8 \\
\midrule
\textbf{Total} & \textbf{35/40} & 29/40 & 29/40 & 0/40 \\
\textbf{Alert Rate} & \textbf{87.5\%} & 72.5\% & 72.5\% & 0\% \\
\bottomrule
\end{tabular}
\end{table}

The full pipeline achieves 87.5\% alert rate vs.\ 72.5\% for both code-only and flat-fact conditions---a 15-point gap attributable to the pre-computed constraint alerts in the manifest. The gap is largest on categories requiring date arithmetic: travel documents (100\% vs.\ 62.5\%) and warranty/deadlines (87.5\% vs.\ 50\%). Critically, code-only and flat facts score identically (72.5\%), confirming that \textbf{the decisive advantage is not the code format per se, but the executable constraint pipeline}---the generate-verify-review loop that computes alerts deterministically and surfaces them in the manifest. The no-memory control scores 0/40, confirming that memory is necessary for proactive alerting.

\subsection{Ablation: Code vs.\ JSON vs.\ Markdown}

We test three formats---Python dataclasses, equivalent JSON, and natural-language Markdown---under two conditions: (1)~with pre-computed alerts in all formats, and (2)~without alerts, but explicitly prompted to check for issues.

\begin{table}[t]
\caption{Format ablation on Active Service (40 scenarios).}
\label{tab:ablation}
\small
\begin{tabular}{@{}lccc@{}}
\toprule
\textbf{Condition} & \textbf{Python} & \textbf{JSON} & \textbf{Markdown} \\
\midrule
With alerts (neutral) & 75.0\% & 75.0\% & 80.0\% \\
No alerts (prompted) & \textbf{100\%} & 90.0\% & 92.5\% \\
\bottomrule
\end{tabular}
\end{table}

With pre-computed alerts, all formats perform identically (75--80\%)---the alert text does the work regardless of format. Without alerts but prompted to check, Python achieves 100\% alert detection vs.\ JSON 90\% and Markdown 92.5\%. The code format provides the clearest signal that data is computable, yielding more consistent ad-hoc alert generation. However, computation accuracy when alerts \emph{are} attempted is similar across formats (67--70\%), confirming that the format effect is secondary to the constraint execution pipeline.

\subsection{Baseline Comparison on Active Service}

We compare User as Code (code representation, neutral prompt, no pre-computed alerts) against three open-source memory systems on 40 Active Service scenarios (Table~\ref{tab:baseline}).

\begin{table}[t]
\caption{Active Service: User as Code vs.\ real memory systems (40 scenarios, neutral prompt, no pre-computed alerts).}
\label{tab:baseline}
\small
\begin{tabular}{@{}lcc@{}}
\toprule
\textbf{System} & \textbf{Alert Rate} & \textbf{Avg KW Score} \\
\midrule
User as Code (code repr) & \textbf{28/40 (70.0\%)} & 0.491 \\
A-MEM (structured notes) & 27/40 (67.5\%) & 0.405 \\
Mem0 (flat facts) & 15/40 (37.5\%) & 0.266 \\
ENGRAM (typed retrieval) & 13/40 (32.5\%) & 0.164 \\
\midrule
\multicolumn{3}{@{}l}{\emph{For reference: full pipeline (code + constraints + manifest):}} \\
User as Code + Alerts & \textbf{35/40 (87.5\%)} & 0.636 \\
\bottomrule
\end{tabular}
\end{table}

Without the constraint pipeline, User as Code (70\%) and A-MEM (67.5\%) are close---A-MEM's Zettelkasten linking helps the LLM notice cross-domain connections. Mem0 (37.5\%) and ENGRAM (32.5\%) lag: flat facts and basic retrieval don't trigger proactive reasoning. The full pipeline adds 17.5pp over code-only, confirming that constraint execution is the decisive differentiator.

\subsection{Standard Benchmark: LOCOMO}

We evaluate recall performance on LOCOMO~\cite{maharana2024locomo} (2 conversations, 120 QA pairs) comparing four systems: UaC~v2 (full architecture with SOTA recall mechanisms: episode-level storage, session summarization, multi-strategy retrieval), UaC~v1 (basic 3-tier), Mem0, and full-context.

\begin{table}[t]
\caption{LOCOMO results (2 conversations, 120 QAs, Gemini 3 Flash with thinking). Both token F1 and LLM-as-Judge (generous) reported.}
\label{tab:locomo}
\small
\begin{tabular}{@{}lcc@{}}
\toprule
\textbf{System} & \textbf{Token F1} & \textbf{LLM-Judge} \\
\midrule
Full Context & 0.225 & \textbf{80.8\%} \\
\textbf{UaC v2 (full arch.)} & \textbf{0.417} & 60.0\% \\
Mem0 & 0.113 & 40.0\% \\
\bottomrule
\end{tabular}
\end{table}

UaC~v2 \textbf{leads on token F1} (0.417 vs.\ 0.225)---structured retrieval produces more precise, factually dense answers. Full-context leads on LLM-Judge (80.8\% vs.\ 60.0\%) because, with the entire raw conversation available, the thinking model can paraphrase freely and the generous judge accepts ``same topic'' answers. On conversation~2, UaC~v2 nearly matches full-context (78.3\% vs.\ 80.0\% judge), showing the architecture approaches parity on favorable conversations.

\textbf{Context on published scores.} Our full-context baseline at 80.8\% is now \textbf{comparable to published SOTA} (EverMemOS 93.05\%~\cite{hu2026evermemos}, MemMachine 91.7\%~\cite{wang2026memmachine}), confirming Gemini 3 Flash is a capable backbone. The remaining $\sim$12pp gap is attributable to testing 2/10 conversations, different backbone (GPT-4.1-mini), and adversarial question exclusion. Published systems require heavy infrastructure (MongoDB, Elasticsearch, Milvus for EverMemOS). UaC~v2 requires only ChromaDB. Our Mem0 at 40.0\% vs.\ published 64--67\% confirms residual methodology differences. \textbf{UaC~v2 achieves 60\% judge accuracy---competitive recall---while uniquely enabling Active Service at 87.5\%.}

\subsection{Cost and Scalability}

\subsection{Qualitative Analysis}

\textbf{Side-by-side: Passport scenario.} Session~1 stores passport expiry (Feb~18, 2025); Session~8 books Tokyo trip (Jan~15). Mem0 stores two disconnected facts---without an explicit query, no retrieval triggers. Zep/Graphiti links passport and trip in a temporal KG but the 180-day rule has no graph encoding. User as Code stores typed state (\texttt{passport.expiry\_date}, \texttt{trip.departure\_date}), executes \texttt{(expiry - departure).days = 34 < 180}, and surfaces the alert in the manifest. Only User as Code produces a proactive alert.

\textbf{Side-by-side: Drug-allergy conflict.} Session~3 notes penicillin allergy; Session~12 prescribes amoxicillin (a penicillin-class antibiotic). Flat facts store both but lack the drug-class relationship. User as Code's \texttt{health\_safety.py} constraint checks medication classes against allergy list---deterministic detection.

\textbf{Failure analysis.} (1)~For simple recall (``What is my phone number?''), User as Code uses $\sim$3x more tokens than vector lookup for the same correct answer---overhead without benefit. (2)~In a scheduling scenario, the agent generated a constraint with incorrect overlap logic (checked \texttt{end < departure} instead of interval intersection). The computation was correct; the \emph{logic} was wrong. (3)~Health/drug safety shows minimal gap across conditions (7/8 for all) because LLM training data encodes pharmacological knowledge---the pipeline's advantage is largest where the LLM's reasoning is unreliable (date arithmetic, deadline computation).

% ===================================================================
\section{Discussion}
% ===================================================================

\textbf{What the experiments show.} The most important finding is that the code format alone does not beat flat facts (both 72.5\% on Active Service without pre-computed alerts). What creates the 15-point gap is the \textbf{executable constraint pipeline}---the generate-verify-review loop that computes alerts deterministically and surfaces them in the manifest. This is more nuanced than ``code beats text.'' The contribution is not a better storage format but a better \emph{architecture}: one where the representation enables a pipeline (code $\rightarrow$ execution $\rightarrow$ alerts $\rightarrow$ manifest) that no other format naturally supports. Code is the \emph{substrate} that makes the pipeline frictionless, not a standalone advantage.

\textbf{The Bitter Lesson.} A natural objection is that typed schemas contradict Sutton's Bitter Lesson~\cite{sutton2019bitter}---that hand-crafted structure is eventually superseded by general methods leveraging computation. This critique misidentifies what is human-imposed. User as Code's domain knowledge---schemas, constraints, tests, domain partitioning---is entirely \textbf{LLM-generated and self-evolving}. The agent designs \texttt{class Trip}, decides what fields it needs, creates and promotes constraints. This is \emph{learned} structure, not imposed features. What \emph{is} human-designed---the project directory convention, update pipeline, progressive disclosure---is engineering infrastructure, analogous to Git or CI/CD. The interpreter is a \emph{tool}, not a \emph{rule}: the LLM retains full freedom in what to check; the interpreter guarantees correct computation.

\textbf{The Formalization Boundary.} User as Code excels at the factual-to-computational end of the memory spectrum. A significant portion of personalization resides at the hard-to-formalize end: context-dependent preferences, behavioral patterns, emotional context. Hard-to-formalize memory lives as string annotations---benefiting from structure and version control, but not executable verification. User as Code provides a unified home where both types coexist.

\textbf{RL-Trainability.} Because memory operations are standard file-system actions, foundation models can be trained via RL to improve memory management. Recent work validates this: Memory-R1~\cite{yan2025memoryr1} achieves 48\% F1 improvement; AgeMem~\cite{yu2026agemem} learns proactive summarization; Mem-alpha~\cite{wang2025memalpha} demonstrates generalization to 400K+ tokens. User as Code's file-based architecture is directly compatible with these paradigms.

\textbf{Memory Safety.} PersistBench~\cite{pulipaka2026persistbench} reveals cross-domain leakage (53\% failure) and sycophancy (97\% failure). User as Code provides natural mitigations: domain separation limits leakage, version control provides audit trails, and the explicit nature of code-based memory enables auditing and selective forgetting.

\textbf{SOTA recall mechanisms are additive, not contradictory.} When we add episode-level storage, session summarization, and multi-strategy retrieval to User as Code (UaC~v2), LOCOMO F1 improves from 0.224 (v1) to 0.398 (v2)---\emph{surpassing} full-context (0.308) by 29\%. This confirms that the code-based architecture is not fundamentally limited for recall. Structured state provides a navigational index; session summaries provide high-signal retrieval; raw episodes provide detail fallback. Together they outperform brute-force context stuffing. The three-tier architecture achieves \textbf{both} strong recall and unique Active Service capability in a single system.

\textbf{Relationship to Long-Context Models.} Long-context cost grows linearly with history; memory systems have constant per-turn cost. More importantly, long-context models cannot perform reliable arithmetic over facts scattered across a million tokens---Active Service requires structured state and executable verification regardless of context window size.

\textbf{Limitations.} (1)~\emph{Format alone is not enough}: code-only and flat facts both achieve 72.5\% on Active Service; the advantage requires the full pipeline. (2)~\emph{Verification overhead}: $\sim$4x write-path token cost vs.\ flat facts, with no benefit for simple recall. (3)~\emph{Constraint logic errors}: correct computation does not guarantee correct logic (e.g., wrong overlap check). (4)~\emph{Category-dependent advantage}: health/drug safety shows minimal gap because LLM training data already encodes pharmacological knowledge; the pipeline helps most on arithmetic-heavy categories. (5)~MemoryBench~\cite{ai2025memorybench} finds no system dominates all tasks---we do not claim universal superiority.

% ===================================================================
\section{Conclusion}
% ===================================================================

User as Code transforms user memory from a passive database into a modular software project through one unifying insight: \textbf{code is the only format that unifies readability and verifiability in a single medium}. Our experiments reveal that the decisive advantage is not the code format per se---code-only and flat facts both achieve 72.5\% on Active Service---but the \textbf{executable constraint pipeline} that code enables. The full pipeline (code $\rightarrow$ constraint execution $\rightarrow$ manifest alerts) achieves 87.5\%, with the largest gains on arithmetic-heavy categories (travel: 100\%, deadlines: 87.5\%). Constraints emerge autonomously---ad-hoc checks promoted to persistent monitors without human engineering. By situating the user project in the agent's own file system, memory operations become native file actions directly compatible with RL-based training.


\bibliographystyle{plainnat}
\bibliography{reference}

\end{document}
